\scanpage{53}
% archive foliation: 52
\[
\begin{array}{l}
\left\{\begin{array}{l}
\displaystyle\sum_1^i\lambda'_j\leq\displaystyle\sum_1^i\lambda_j\\
\displaystyle\sum_{\unsure{1}}^{\unsure{n}}\lambda'_j=\displaystyle\sum\lambda_j
\end{array}\right.\\
\qquad\unsure{\Vert}\\
\qquad\displaystyle\sum i\unsure{h'_i}
\end{array}
\]
\margin{\unsure{mult de} pentes de Hodge}
\figure{\begin{tikzpicture}\draw (0,0)--(3,.9) .. controls (3.2,1) and (3.1,1.1) .. (3.05,1.1);\end{tikzpicture}}
\hrule

\margin{\unsure{Sur un th de Tate.}}
\uline{Sur un théorème de Tate.} (Devrait venir avec foncteur mystérieux !)

C'est le théorème suivant (Driebergen ....):

\uline{Th.} (Tate): Soit S un schéma noethérien normal intègre dont le corps des fonctions \inserted{$k(\eta)$} est de car. zero, Alors le foncteur
\[
\mathrm{BT}(S)\longrightarrow\mathrm{BT}(\eta)
\]
est pleinement fidèle.

Ainsi, un groupe de BT sur S est connu \inserted{à isom unique près} quand on connaît la représentation p-adique qu'\deleted{elle}\inserted{\unsure{il}} définit de $\pi_1(\eta)=\mathrm{Gal}(\overline{k(\eta)}/k(\eta))$. Prenant d'abord le cas où S est le spectre d'un anneau de valuation discrète complet V d'inégales car. , de corps \deleted{des} résiduel k parfait, la question se pose donc

a) De caractériser les modules galoisiens qu'on peut obtenir (Tate donne des conditions nécessaires: les modules en question doivent être du type de ``Hodge-Tate'' );

b) De trouver un foncteur ``explicite'' sur la catégorie de ces modules galoisiens, vers les F-cristaux sur k M muni d'une filtration de $M\otimes_W\unsure{W|V}$.

\margin{\uline{Problème}}
Le théorème de Tate suggère de plus que le ``foncteur mystérieux'' qui va des F-cristaux sur k filtrés sur L, vers les $\underline{Q}_p$-modules galoisiens sur L, est pleinement fidèle, et qu'on doit pouvoir donc \deleted{définir un \unsure{foncteu|foncteur}} expliciter un foncteur quasi-inverse (non moins mystérieux), allant d'une sous-catégorie pleine convenable des $\underline{Q}_p$-modules galoisiens sur L (sans doute contenue dans celle des modules galoisiens du type de Hodge-Tate) vers les F-cristaux sur k filtrés sur L.

D'autre part, on conjecture que dans le th. de Tate ci-dessus, la restriction au cas où $\deleted{k}\inserted{k(\eta)=L}$ est de car. nulle est inutile, Le seul cas qui offre

